\section{Solving the next growth interval with a diffusing parent}
\label{sec:cv-diffusive-parent}

This section continues an uncompensated equal-diffusion variant of the
first stage through the next growth interval. It keeps the laboratory
direction and the original coefficients. The exact amplitude system
is realized below both on a forced affine parent and on an actual
two-wave state with an explicitly computed interaction force.
The latter force specifies the difference between the affine-gradient
calculation and the full spatial equations.

\subsection{The inherited gradient and the original second pair}

Retain the preceding section's $A_0,\mu,\nu,K_1,s,\lambda_1$, and
let $\nu_{\rm phys}=\kappa_{\rm phys}=d>0$.
For a sine first wave on its rotating affine base, the exact amplitude
equations with retained diffusion are
\[
 \dot\Theta_{1,d}=a_1\Omega_{1,d}-dK_1^2\Theta_{1,d},\qquad
 \dot\Omega_{1,d}=b_1\Theta_{1,d}-dK_1^2\Omega_{1,d}.
\]
If $(\Theta_1,\Omega_1)$ denotes the first inviscid pair of the
preceding section on its specified control and times, direct
differentiation proves
\[
 (\Theta_{1,d},\Omega_{1,d})(t)
 =e^{-dK_1^2(t-t_0)}(\Theta_1,\Omega_1)(t).
\]
Both the forward map and its inverse multiplication by the positive
exponential satisfy the indicated systems, so they are a bijection
of solution spaces. In particular the selected vorticity still returns
to zero at the specified $t_b$. This uses the inviscid times and does
not assert the damped temperature crosses the same growth threshold.
Its exact log-gain from the original seed is
\[
 \log(P_b/S)-\frac{dK_1^2}{\nu\sin s}
                       (L+2+\Lambda_1^{-1}).
\]

Set $P_{d,b}=e^{-dK_1^2(t_b-t_0)}P_b>0$,
$A_{1,d}=K_1P_{d,b}$, and $r=t-t_b\geq0$.
During the fixed first holding control $\alpha=s$ the exact first
pair is
\[
 \Theta_{1,d}(r)=-P_{d,b}e^{-qr},\quad\Omega_{1,d}=0,
 \quad q=dK_1^2,
 \quad G_d(r)=(A_0\sin s,-A_0\cos s-A_{1,d}e^{-qr}).
\]
The gradient therefore changes at the rate
$G_d'(r)=(0,qA_{1,d}e^{-qr})$.
Let the new laboratory phase be $\zeta_2=e(s_2)$ with
$0<s_2\leq1/8$, and let $K_2=\mu\lambda_2>0$.
The actual first hold has zero velocity gradient, so $\zeta_2'=0$.
The second affine-wave coefficients are exactly
\begin{equation}
 a_{2,d}(r)=-\frac{J\zeta_2\cdot G_d(r)}{K_2}
 =\frac{A_{1,d}e^{-qr}\sin s_2+A_0\sin(s+s_2)}{K_2},\quad
 b_2=K_2\sin s_2,\quad\delta_2=dK_2^2.
 \label{eq:cv-feedback-coefficients}
\end{equation}
Their signs are positive because $0<s+s_2\leq1/4$.
The exact second sine amplitude system is
\begin{equation}
 \Theta_2'=a_{2,d}(r)\Omega_2-\delta_2\Theta_2,
 \qquad\Omega_2'=b_2\Theta_2-\delta_2\Omega_2.
 \label{eq:cv-feedback-pair}
\end{equation}
In particular the changing inherited gradient has not been replaced
by its endpoint value. An affine realization alone would require the
scalar force $G_d'(r)\cdot x$: the affine Laplacian is zero, so that
force cannot be inferred to vanish from the old sine-wave diffusion.

\subsection{Exact operator change and convergent fundamental solutions}

Define
\[
 c_0=A_0\sin(s+s_2)\sin s_2>0,\qquad
 c_1=A_{1,d}\sin^2s_2>0,\qquad W(r)=e^{\delta_2r}\Omega_2(r).
\]
Because $b_2$ is constant, differentiating the second original
equation gives the exact system equivalence
\begin{equation}
 W''=(c_0+c_1e^{-qr})W,\qquad
 \Omega_2=e^{-\delta_2r}W,\qquad
 \Theta_2=\frac{e^{-\delta_2r}}{b_2}W'.
 \label{eq:cv-feedback-scalar-map}
\end{equation}
Conversely these last two formulas give
$\Omega_2'=b_2\Theta_2-\delta_2\Omega_2$ and
$\Theta_2'=e^{-\delta_2r}W''/b_2-\delta_2\Theta_2
=a_{2,d}\Omega_2-\delta_2\Theta_2$, since
$b_2a_{2,d}=c_0+c_1e^{-qr}$.
Thus the initial-data map is the invertible map
$(\Theta_2(0),\Omega_2(0))\mapsto(W(0),W'(0))
=(\Omega_2(0),b_2\Theta_2(0))$.

Put
\[
 x(r)=x_0e^{-qr/2},\quad x_0=\frac{2\sqrt{c_1}}q,\quad
 \vartheta=\frac{2\sqrt{c_0}}q>0,\qquad W(r)=\mathcal W(x(r)).
\]
Here $x(r)$ is a scalar change of independent variable and is unrelated
to either physical spatial coordinate $x_1,x_2$, which remain fixed
in the PDE below. Since $x'=-qx/2$ and $x''=q^2x/4$, one has
\[
 W''-(c_0+c_1e^{-qr})W
 =\frac{q^2}4\left[x^2\mathcal W_{xx}+x\mathcal W_x
                         -(x^2+\vartheta^2)\mathcal W\right].
\]
This proves the operator transformation including its nonzero factor.
Rather than invoking a special-function name, define on $x>0$
\begin{equation}
 f(x)=x^{\vartheta}\sum_{n=0}^\infty a_nx^{2n},\quad
 a_0=1,\quad a_n=\frac{a_{n-1}}{4n(n+\vartheta)},\qquad
 g(x)=f(x)\int_{x_0}^x\frac{d\xi}{\xi f(\xi)^2}.
 \label{eq:cv-feedback-fundamental}
\end{equation}
For every compact interval in $x>0$, the consecutive-term absolute
ratio $x^2/[4n(n+\vartheta)]$ tends uniformly to zero.
Every fixed derivative multiplies these terms by a polynomial in $n$
and a bounded power of $x^{-1}$ on that compact interval; its series
therefore converges uniformly as well. Termwise substitution is valid.
The coefficient of $x^{\vartheta+2n}$ in
$x^2f''+xf'-(x^2+\vartheta^2)f$ is
$4n(n+\vartheta)a_n-a_{n-1}=0$ for $n\geq1$; the $n=0$
coefficient is also zero. All coefficients are positive, so $f>0$.
The defining integral for $g$ consequently exists on every compact
subinterval of $x>0$. If $I'=1/(xf^2)$, differentiation yields
\[
 f(fI)' - f'(fI)=f^2I'=1/x.
\]
Also $I''/I'=-1/x-2f'/f$; substituting $g=fI$ into
$g''+g'/x-(1+\vartheta^2/x^2)g$ leaves
$fI''+(2f'+f/x)I'=0$. Thus $g$ solves the same equation and its
Wronskian with $f$ is the nonvanishing function $1/x$.
For completeness this proves completeness of the two solutions:
at any $x>0$ their value-and-derivative matrix is invertible, so one
linear combination matches arbitrary Cauchy data; linear uniqueness
then identifies that combination with any solution on overlapping
compact subintervals of $(0,\infty)$.

For the original data let $W_0=\Omega_2(0)$ and
$W_1=b_2\Theta_2(0)$. The exact solution is
\begin{equation}
 W(r)=c_ff(x(r))+c_gg(x(r)),\quad
 c_f=\frac{W_0}{f(x_0)},\quad
 c_g=-\frac{2f(x_0)}q\left(W_1+\frac{qx_0}{2}c_ff'(x_0)\right).
 \label{eq:cv-feedback-initial-map}
\end{equation}
Indeed $g(x_0)=0$, $g'(x_0)=1/[x_0f(x_0)]$, and
$x'(0)=-qx_0/2$ verify both initial conditions directly.
For every finite $r$, $x(r)>0$, so the displayed functions and their
derivatives are finite there. If $d=0$, the preceding substitutions
dividing by $q$ are not used: the original equation is
$W''=(c_0+c_1)W$, with the exact solution
$W=W_0\cosh(\sqrt{c_0+c_1}\,r)
+W_1\sinh(\sqrt{c_0+c_1}\,r)/\sqrt{c_0+c_1}$.

\subsection{Signs and quantitative physical gain with this feedback}

Let $\gamma_0=\sqrt{c_0}$ and $\gamma_i=\sqrt{c_0+c_1}$.
Use the original instantaneous growing eigenline at entry:
$\Theta_2(0)<0$ and
$\Omega_2(0)=\sqrt{b_2/a_{2,d}(0)}\Theta_2(0)<0$.
Then $Y=-W$ has $Y(0)=Y_0>0$ and
$Y'(0)=\gamma_iY_0>0$.
As long as $Y>0$, its equation gives $Y''>0$ and
$Y'\geq Y'(0)>0$; this excludes a first zero by continuity.
Define
\[
 Y_-=Y_0\left(\cosh(\gamma_0r)
          +\frac{\gamma_i}{\gamma_0}\sinh(\gamma_0r)\right),
 \qquad Y_+=Y_0e^{\gamma_i r}.
\]
The differences have zero initial value and derivative and satisfy
\begin{align*}
 (Y-Y_-)''-\gamma_0^2(Y-Y_-)&=c_1e^{-qr}Y,\\
 (Y_+-Y)''-\gamma_i^2(Y_+-Y)&=c_1(1-e^{-qr})Y.
\end{align*}
For $a>0$, the formula
$Z(r)=\int_0^r\sinh(a(r-u))S(u)/a\,du$ has
$Z(0)=Z'(0)=0$ and $Z''-a^2Z=S$, by two differentiations.
Uniqueness proves it is the solution of that initial-value problem.
Both sources above are nonnegative; both this kernel and its first
$r$ derivative are nonnegative for $0\leq u\leq r$.
Therefore $Y_-\leq Y\leq Y_+$ and
$Y_-'\leq Y'\leq Y_+'$ for every finite $r\geq0$.
The inverse physical map now gives both amplitudes negative and the
precise temperature-gain bounds
\begin{equation}
 e^{-\delta_2r}\left(\cosh(\gamma_0r)
       +\frac{\gamma_0}{\gamma_i}\sinh(\gamma_0r)\right)
 \leq\frac{-\Theta_2(r)}{-\Theta_2(0)}
 \leq e^{(\gamma_i-\delta_2)r}.
 \label{eq:cv-feedback-physical-gain}
\end{equation}
In particular, if the actual numerical coefficients satisfy
$\delta_2\geq\gamma_i$, this particular growth control never raises
the second temperature above its entry amplitude. This is a proved
obstruction to that specified uncompensated control, not a statement
about other controls or the existence of an infinite viscous sequence.

\subsection{The exact map to spatial fields and its interaction force}

The affine parent $\theta_b=G_d(r)\cdot x$, $u_b=0$ is a solution
with pressure $p_b=x_2G_d(r)\cdot x/2$ and forces
$f_{\theta,b}=G_d'(r)\cdot x$ and
$f_{u,b}=-(A_0\sin s/2)Jx$.
The scalar statement follows from $\Delta\theta_b=0$.
For momentum, direct differentiation gives
$\nabla p_b-\theta_be_2=(G_dx_2-e_2(G_d\cdot x))/2
=-G_{d,1}Jx/2$.
An added sine wave whose amplitudes solve
\eqref{eq:cv-feedback-pair} has zero additional PDE residual except
for its activation, by the same explicit phase, transverse-vector,
and pressure calculations below. Thus the solved system has an exact
forced affine-background realization.

There is also a full two-wave realization specifying the original
spatial discrepancy. Let $F_1,F_2$ be smooth odd periodic profiles,
each with mean-zero primitive $Q_j'=F_j$, and assume $F_1'(0)=1$.
This includes both sine and the released profile which is linear
near zero. Set
\[
 G_0=(A_0\sin s,-A_0\cos s),\quad
 T(r)=-P_{d,b}e^{-qr},\quad
 q_1=K_1x_2,\quad q_2=K_2\zeta_2\cdot x,
\]
and let $w_2(r)$ be any specified smooth activation, for example
the source's $\eta(\Gamma_2r)$. Define on $\mathbb R^2$
\begin{align}
 \theta&=G_0\cdot x+T F_1(q_1)+w_2\Theta_2F_2(q_2),\nonumber\\
 u&=\frac{w_2\Omega_2}{K_2}J\zeta_2F_2(q_2),\nonumber\\
 p&=\frac12x_2(G_0\cdot x)+\frac{T}{K_1}Q_1(q_1)
          +\frac{w_2\Theta_2\zeta_{2,2}}{K_2}Q_2(q_2).
 \label{eq:cv-feedback-spatial-state}
\end{align}
Their complete retained-diffusion forces are
\begin{align}
 f_\theta={}&-dK_1^2T\,[F_1(q_1)+F_1''(q_1)]
 +w_2'\Theta_2F_2(q_2)\nonumber\\
 &-dK_2^2w_2\Theta_2\,[F_2(q_2)+F_2''(q_2)]
 +f_{\rm cross},\label{eq:cv-feedback-spatial-scalar-force}\\
 f_{\rm cross}={}&\frac{w_2\Omega_2}{K_2}\sin s_2\,
        K_1T\,[F_1'(q_1)-1]F_2(q_2),
 \label{eq:cv-feedback-exact-cross-force}\\
 f_u={}&-\frac{A_0\sin s}2Jx
 +\frac{w_2'\Omega_2}{K_2}J\zeta_2F_2(q_2)
 -dK_2w_2\Omega_2J\zeta_2[F_2(q_2)+F_2''(q_2)].
 \label{eq:cv-feedback-spatial-momentum-force}
\end{align}
Here is a direct verification. The old wave has derivative $T'=-dK_1^2T$
and Laplacian $K_1^2TF_1''$, giving the first scalar force.
The old temperature gradient differs from the coefficient gradient
$G_d=G_0+K_1Te_2$ by precisely
$K_1T[F_1'(q_1)-1]e_2$. Its advection by $u$ is the displayed
$f_{\rm cross}$ because $J\zeta_2\cdot e_2=\sin s_2$.
The remainder of that old-temperature advection is
$(w_2\Omega_2/K_2)J\zeta_2\cdot G_d\,F_2=-w_2a_{2,d}\Omega_2F_2$;
it cancels the coupling part of $w_2\Theta_2'F_2$.
The new wave advects its own temperature and velocity by zero since
$J\zeta_2\cdot\zeta_2=0$. Its amplitude damping and its Laplacian
give the two new-profile terms in the forces.
The old wave pressure has gradient $TF_1(q_1)e_2$, so it cancels
old wave buoyancy exactly. The new pressure has gradient
$w_2\Theta_2\zeta_{2,2}\zeta_2F_2$.
The new momentum coupling vector before that pressure is
$w_2\Theta_2(\zeta_{2,1}J\zeta_2-e_2)F_2
=-w_2\Theta_2\zeta_{2,2}\zeta_2F_2$, by resolving $e_2$ in the
orthonormal basis $(\zeta_2,J\zeta_2)$. Hence it cancels as well.
The base pressure supplies exactly the first vector-force term,
as proved above, and divergence is zero by transversality.
This verifies both full equations including all interactions.

For the sine profiles, $F_j+F_j''=0$ identically. The remaining
scalar interaction force is then exactly
\[
 f_{\rm cross}=-\frac{A_{1,d}e^{-qr}}{K_2}
   w_2\Omega_2\sin s_2\,[\cos(K_1x_2)-1]\sin(q_2),
 \qquad
 \|f_{\rm cross}(r)\|_\infty
 \leq\frac{2A_{1,d}e^{-qr}}{K_2}|w_2\Omega_2|\sin s_2.
\]
The bound follows from $|\cos u-1|\leq2$ and $|\sin v|\leq1$.
For the released linear-core profile the same explicit formulas hold,
and $F_1'-1=0$ on its original linear interval, so the interaction
force vanishes there. Its old-profile diffusion defect is instead
$-dK_1^2Tq_1$ on that interval, since $F_1(q_1)=q_1$ and
$F_1''(q_1)=0$. Its gradient is exactly $G_d'(r)$.
Thus exponential damping of that released linear core requires the
stated core force, and is not an unforced consequence of its zero
Laplacian. The explicit forces, inverse solution maps, and convergent
formulas above carry the feedback calculation to an actual spatial
solution without conflating the sine state, its affine gradient, and
the released profile. No compact-support or infinite-iteration claim
is made in this section.
